\documentclass[11pt,a4paper]{article}
\usepackage[margin=2.6cm]{geometry}
\usepackage{amsmath,amssymb}
\usepackage[authoryear,round]{natbib}
\usepackage{xcolor}
\usepackage[colorlinks=true,linkcolor=blue!50!black,citecolor=blue!45!black]{hyperref}

\newcommand{\PB}{P^{\mathrm{B}}}
\newcommand{\PJ}{P^{\mathrm{J}}}
\newcommand{\assoc}{\operatorname{assoc}}
\newcommand{\bigO}{O}

\title{Positioning in Economics\\
\large Companion to \emph{The Paper's Dialectic}, for
\emph{Order-Sensitivity of Belief and Decision Statistics under Jeffrey
Conditioning}}
\author{}
\date{}

\begin{document}
\maketitle
\vspace{-2.4em}

\begin{center}\small\itshape
Standalone working document, not manuscript content. The philosophical
standoff---D\"oring's attack, Hawthorne's refusal, Wagner's defence---is treated
in the companion note and is not repeated here. Every paper cited below was read
directly; the two whose results are formalized carry a pointer to the
formalization.
\end{center}

\section*{Why economics is the second audience}

The companion note leaves the debate over amnestic updating in a particular
state. Its participants argue about realism, and none of them measures anything.
The paper answers the answerable half of that debate by computing where a
sequence effect shows up in a population's statistics and where it cannot. That
is a claim about what an instrument can detect, and questions of that form have
a settled home in economics, where three literatures have independently arrived
at the same discipline.

\paragraph{Non-identification is a recognised result, not an excuse.}
\citet{ShmayaYariv2016} show that absent restrictions on subjects' conjectures
about the experimental design, \emph{every} mapping from signal sequences to
chosen alternatives can be rationalised by Bayesian updating---their Theorem~2,
the ``anything goes'' result---while a restriction on the design restores
testable implications through a no-reversal condition (their Theorem~1).
\citet{Heckman1998} showed that audit-pair estimates of discrimination rest on an
assumption about unobserved productivity that ``nothing guarantees,'' and that
under threshold hiring the audit ``can find discrimination when in fact none
exists; it can also disguise discrimination when it is present'' (p.~102).
\citet{Bohren2019} observe that in static settings ``different sources generate
the same patterns of observable behavior,'' so the source of discrimination is
not identified from such data at all. In each case the obstruction is structural:
no sample size repairs it.

\paragraph{The present paper's obstruction is of that kind.}
The believed cross-attribute association departs from its sequence-free value
only at $\bigO(c^{2})$, while the marginals and the decision rate depart at
$\Theta(c)$. The first-order signal in the protected statistics is zero in the
population itself, so the failure is one of identification and not of estimation.
An audit that reads believed association and finds conformity has learned
nothing about whether its evaluators update sequentially, and would learn nothing
from a hundred times as many evaluators. What differs from the three results
above is the source of the obstruction. Heckman's turns on unobserved
heterogeneity meeting a threshold, and is repaired by assumptions;
Shmaya--Yariv's turns on unobserved conjectures about the design, and is repaired
by restricting the protocol. The obstruction here is geometric---the statistic's
differential annihilates the plane the two reading sequences span---and no
auxiliary assumption repairs it. It is removed only by changing the statistic.

\section*{Instrument choice as a genre}

\citet{AugenblickRabin2021} test Bayesian updating through the martingale
property, which for a belief stream gives $E[f(\theta_0,\dots,\theta_t)(\theta_t
-\theta_{t+1})]=0$ \emph{for any instrument} $f$. Different tests arise from
different instruments; theirs amounts to the choice $f=2\theta_t-1$, from which
the equality of expected movement and expected uncertainty reduction follows
(their Proposition~1). They are explicit that they do not claim their instrument
is universally more powerful, and they do not attempt to say which instruments
have power against which deviations.

That is precisely the question the present paper closes for its own deviation.
Proposition~PRO characterises the smooth statistics that fail to register the
sequence effect at first order: exactly those whose differential at independence
points along the differential of the believed association. The condition is
necessary as well as sufficient, so the classification is complete rather than
selective. Two further contrasts are worth stating. First, the test here is
cross-sectional: it needs one measurement per evaluator, with reading order
varying across evaluators, where Augenblick--Rabin need a belief stream over time
for each agent. Second, benchmark-freeness is \emph{not} a distinguishing
feature---their driving assumption, held throughout, is that the
normatively-correct beliefs are unobserved---but it is evidently a property this
literature values, and Proposition~ORD has it.

Augenblick--Rabin also supply the closest published analogue to the protected
class. For streams of \emph{expectations} rather than binary beliefs they prove
that any observed stream is arbitrarily likely under some data-generating
process, so an expectations stream carries no testable content; their diagnosis
is that without a scale for movements there is nothing to test. That is the same
shape of finding as the one here---one summary of a belief retains testable
content about the updating process and another does not---with one instructive
difference. Their negative result is absolute, since no scale exists. The result
here is graded: the association does carry the effect, at an order of magnitude
below the levels, and the prior covariance $c$ supplies exactly the scale their
expectations lack. It is that scale which makes the comparison
$c^{2}\ll\varepsilon\ll c$ meaningful, and which lets every statistic be placed
at an order rather than merely inside or outside the testable class.

\section*{The identifying design, and its precedent}

\citet{Bohren2019} face a non-identification and resolve it by conditioning on a
sequence variable. Static data cannot separate correct-belief, biased-belief and
preference-based discrimination; but along a history of evaluations the three
sources diverge, since correct-belief discrimination attenuates and never
reverses, so an observed reversal identifies bias. Their field experiment
supplies the reversal.

The design proposed here has the same structure. A pooled aggregate cannot
separate sequential from sequence-free updating, but evaluators split by which
cue they met first differ at first order in levels and in decisions while
agreeing to second order in association. The conditioning variable is the
evaluator's reading order, as theirs is the candidate's evaluation history, and
in both cases it is a variable that pooled data discards. The practical
implication is correspondingly cheap: reading order must be recorded, or
randomised, for the question to be answerable at all.

\section*{A rival mechanism, and why the signature survives it}

Order effects in belief updating need not come from the input type. In
\citet{Wilson2014} an agent with finitely many memory states updates optimally,
the extreme states become nearly absorbing, and order effects follow: a block of
identical signals moves final beliefs more when it comes first in short
sequences, and more when it comes last in long ones (her Theorem~4). Since a
Bayesian with unbounded memory treats exchangeable evidence symmetrically, this
is a genuine competing explanation for the observable this paper points at, and
it is rational rather than a bias.

It does not, however, mimic the signature. Wilson's agent carries a single belief
about one state of the world, so her model makes no prediction about the
association between two attributes; the quantity that does the identifying work
here does not exist in it. The prediction under test is a \emph{pair} of orders:
$\Theta(c)$ between reading orders in levels and decisions, together with
$\bigO(c^{2})$ in the believed association. A one-dimensional account of order
effects cannot produce the second half, and an account extended to two attributes
would have to protect the association for its own reasons. The two-part signature
is therefore not decoration; it is what separates this mechanism from its
rivals, and it is the reason the elicitation must include the association even
though the association is the statistic that shows nothing.

\bigskip
\bibliographystyle{plainnat}
\begin{thebibliography}{9}
\bibitem[Augenblick \& Rabin(2021)]{AugenblickRabin2021}Augenblick, N. \& Rabin,
  M. (2021). Belief movement, uncertainty reduction, and rational updating.
  \emph{Quarterly Journal of Economics} 136(2), 933--985. Read: abstract,
  introduction, definitions, and the instrument discussion of \S4; Proposition~1,
  Corollary~1 and Table~1 formalized and reproduced in
  \texttt{literature/augenblick\_rabin2021/}.
\bibitem[Bohren, Imas \& Rosenberg(2019)]{Bohren2019}Bohren, J.~A., Imas, A. \&
  Rosenberg, M. (2019). The dynamics of discrimination: theory and evidence.
  \emph{American Economic Review} 109(10), 3395--3436. Read: abstract and
  introduction, including the identification argument; model and experimental
  sections skimmed.
\bibitem[Heckman(1998)]{Heckman1998}Heckman, J.~J. (1998). Detecting
  discrimination. \emph{Journal of Economic Perspectives} 12(2), 101--116. Read:
  introduction, the definition and measurement section, and the audit-pair
  critique; quotation from p.~102.
\bibitem[Shmaya \& Yariv(2016)]{ShmayaYariv2016}Shmaya, E. \& Yariv, L. (2016).
  Experiments on decisions under uncertainty: a theoretical framework.
  \emph{American Economic Review} 106(7), 1775--1801. Read: abstract, overview,
  model, and Theorems~1--2; both theorems' cores formalized in
  \texttt{literature/shmaya\_yariv2016/}.
\bibitem[Wilson(2014)]{Wilson2014}Wilson, A. (2014). Bounded memory and biases in
  information processing. \emph{Econometrica} 82(6), 2257--2294. Read: abstract,
  introduction, and the order-effect results of \S4 (Theorem~4); the
  working-paper version in the project folder is dated 2003.
\end{thebibliography}

\end{document}
